% =============================================================
%  Justificativa da Pesquisa (1,00 ponto)
% =============================================================
\section*{Justificativa da Pesquisa (1,00 ponto)}

A justificativa deste projeto parte de demandas relatadas por peritos criminais em atividade, levantadas junto a colegas da área de ciências forenses:

\begin{itemize}[leftmargin=1.5em]
    \item \textbf{A análise é mais custosa que o cálculo.} Segundo relato direto de perito atuante, cálculos como o de velocidade média são tecnicamente simples; o esforço real está em localizar os trechos relevantes de longas gravações, capturar e organizar frames, interpretar o conteúdo visual e montar uma demonstração compreensível para o laudo --- um processo hoje inteiramente manual e repetitivo.
    \item \textbf{Câmeras de monitoramento urbano não são padronizadas para fins periciais.} Diferentes fabricantes, resoluções, ângulos e taxas de quadros aumentam a complexidade da análise em comparação a sistemas dedicados (como radares ou câmeras de trânsito).
    \item \textbf{Alto volume de solicitações.} Exames de vídeo estão entre as perícias de imagem mais requisitadas, especialmente em investigações que dependem de câmeras de segurança como única fonte de prova.
    \item \textbf{Amadurecimento de técnicas de processamento de vídeo.} Algoritmos de detecção/rastreamento de objetos e, mais recentemente, modelos multimodais de linguagem para vídeo já demonstram capacidade de descrever e responder perguntas sobre conteúdo visual, mas sua aplicação sistemática e comparativa ao contexto pericial brasileiro ainda é pouco explorada na literatura.
    \item \textbf{Urgência investigativa.} Reduzir o tempo da etapa de análise impacta diretamente prazos de inquérito e a celeridade da produção da prova.
\end{itemize}

Dessa forma, uma análise comparativa de algoritmos que possam complementar um modelo de linguagem no apoio ao trabalho do perito tem potencial de impacto direto na rotina pericial, contribuindo para a padronização, a celeridade e a acessibilidade técnica desses exames.
